% !TEX root = ../main.tex

\thispagestyle{empty} % убираем колонтитулы с титульного листа

\begin{center}
	\workMinistry \\
	\workUniversity \\
	\workInstitute
\end{center}

\vspace{2ex}

\hfill\begin{minipage}{7.5cm}
	УТВЕРЖДАЮ\\[1ex]
	Директор ВШ КТиИС\\[1ex]
	\underline{\hspace{3cm}} М.\,А.~Полтавцева\\[1ex]
	«\underline{\hspace{1cm}}»\underline{\hspace{3cm}} 2026 г.
\end{minipage}

\vspace{4ex}

\begin{center}
	\textbf{ЗАДАНИЕ}\\
	\textbf{на выполнение выпускной квалификационной работы}
\end{center}

\vspace{2ex}

\noindent
студенту \uline{Плеханову Артёму Евгеньевичу, группа 5130901/20201} \\
\centerline{\scriptsize (фамилия, имя, отчество (при наличии), номер группы)}

\vspace{2ex}

\noindent
1. Тема работы: \uline{Разработка кроссплатформенного мобильного приложения для трекинга достижений с элементами геймификации}

\vspace{1ex}

\noindent
2. Срок сдачи студентом законченной работы: \uline{21.05.2026}

\vspace{1ex}

\noindent
3. Исходные данные по работе: \uline{Требования к функциональности приложения: регистрация и аутентификация пользователя; формирование и редактирование коллекций «паков» достижений; просмотр достижений и шагов выполнения; отметка прогресса с сохранением на сервере; загрузка изображений; локализация интерфейса (русский/английский). Исходные данные к проектированию: спецификация REST API (OpenAPI), описание предметной области «достижения и прогресс», требования к целевым платформам (Android, перспектива Web). Нормативная и учебная литература по многоплатформенной разработке, клиент-серверному взаимодействию и проектированию пользовательских интерфейсов.}

\vspace{1ex}

\noindent
4. Содержание работы (перечень подлежащих разработке вопросов):

\vspace{1ex}

\noindent
\uline{Анализ предметной области и существующих решений для учёта целей и достижений; постановка задачи и формирование требований к системе.}

\vspace{1ex}

\noindent
\uline{Выбор технологий: Kotlin Multiplatform, Compose Multiplatform, клиент HTTP, навигация, хранение настроек и токенов, загрузка изображений.}

\vspace{1ex}

\noindent
\uline{Проектирование архитектуры приложения (слой представления, ViewModel, репозитории, доступ к API).}

\vspace{1ex}

\noindent
\uline{Реализация модулей: аутентификация, работа с коллекциями паков и достижений, экраны списков и деталей, сценарии создания и редактирования.}

\vspace{1ex}

\noindent
\uline{Интеграция с сервером: вызовы API, обработка ошибок, синхронизация прогресса.}

\vspace{1ex}

\noindent
\uline{Реализация пользовательского интерфейса и локализации; обеспечение согласованности состояния при асинхронных операциях.}

\vspace{1ex}

\noindent
\uline{Тестирование и оценка результатов; формулировка выводов и направлений развития (офлайн-режим, расширение платформ и т.п.).}

\vspace{2ex}

\noindent
5. Перечень графического материала (с указанием обязательных чертежей): \uline{Диаграмма компонентов приложения и схема взаимодействия с сервером. Диаграмма основных экранов (навигация, пользовательские сценарии). Скриншоты ключевых экранов работающего приложения. Чертежи по ЕСКД для данной темы не требуются.}

\vspace{1ex}

\noindent
6. Перечень используемых информационных технологий, в том числе программное обеспечение, облачные сервисы, базы данных и прочие сквозные цифровые технологии (при наличии): \uline{Язык программирования Kotlin; Kotlin Multiplatform; Compose Multiplatform; Material Design 3; Gradle (в т.ч. каталог версий зависимостей); Android SDK; среда разработки Android Studio / IntelliJ IDEA; Kotlin Coroutines, Flow, StateFlow; архитектурный подход MVVM, слой репозиториев; Navigation 3 (навигация по экранам); Ktor Client (HTTP); REST, JSON; OpenAPI Generator (клиент по спецификации API); OAuth 2.0, Bearer token (использование на клиенте); Coil (изображения); FileKit (выбор файлов); Multiplatform Settings (локальные настройки); локализация интерфейса (ресурсы строк, многоязычность); система контроля версий Git.}

\vspace{1ex}

\noindent
7. Консультанты по работе: \hrulefill

\vspace{2ex}

\noindent
8. Дата выдачи задания \uline{\hspace{1cm}21.04.2026\hspace{1cm}}

\vspace{2ex}

\noindent
\begin{tabular}{@{}p{0.65\textwidth}l@{}}
Руководитель ВКР & \\
\makebox[5cm]{\hrulefill} & \workTeacher \\
\makebox[5cm]{\centerline{\scriptsize (подпись)}} & \centerline{\scriptsize инициалы, фамилия} \\[2ex]
Задание принял к исполнению \uline{\hspace{1cm}21.04.2026\hspace{4cm}} & \\
\makebox[10cm]{\centerline{\scriptsize (дата)}} & \\[2ex]
Студент & \\
\makebox[5cm]{\hrulefill} & \workAuthors \\
\makebox[5cm]{\centerline{\scriptsize (подпись)}} & \centerline{\scriptsize инициалы, фамилия} \\
\end{tabular}

\newpage